%%%%%%%%%%%%%%%%%%%%%%%%%%%%%%%%%%%%%%%%%%%%%%%%%%%%%%%%%%%%%%%%%%%%%%%%%%%%
%% Text Area: 8in (include Runningheads) x 5in
%% ws-ijcm.tex: 09-12-2024
%% TeX file to use with ws-ijcm.cls written in LaTeX2E.
%% The content, structure, format and layout of this style file is the
%% property of World Scientific Publishing Co. Pte. Ltd.
%% Copyright 2024 by World Scientific Publishing Co.
%% All rights are reserved.
%%%%%%%%%%%%%%%%%%%%%%%%%%%%%%%%%%%%%%%%%%%%%%%%%%%%%%%%%%%%%%%%%%%%%%%%%%%%

%\documentclass[wsdraft]{ws-ijcm}
\documentclass{ws-ijcm}

\usepackage{natbib}
\usepackage{algorithm}
\usepackage{algpseudocode}
\usepackage{xcolor}
\usepackage[verbose,hypertexnames=false]{hyperref}
\hypersetup{colorlinks=false,allbordercolors=blue,pdfborderstyle={/S/U/W 1}}
% \label, \ref and \cite commands are highly recommended

\begin{document}

%%%%%%%%%%%%%%%%%%%%% Publisher's Area please ignore %%%%%%%%%%%%%%%
%
\catchline{}{}{}{}{}
%
%%%%%%%%%%%%%%%%%%%%%%%%%%%%%%%%%%%%%%%%%%%%%%%%%%%%%%%%%%%%%%%%%%%%

\markboth{[Authors]}{PE-PINN: A Physics-Enhanced PINN for PKN Hydraulic Fracture Modeling}

\title{PE-PINN: A Physics-Enhanced Physics-Informed Neural Network \\
for Solving the PKN Model of Hydraulic Fracturing%
\footnote{For the title, try not to use more than 3 lines. Typeset the title in 10 pt bold.}
}

\author{[Author 1]\footnote{Typeset names in 8 pt Roman.}}

\address{[Department, University]\\
[City, Country]\,\footnote{State completely without abbreviations.}\\
\email{[email]}}

\author{[Author 2]}

\address{[Department, University]\\
[City, Country]\\
\email{[email]}}

\maketitle

\begin{history}
\received{(Day Month Year)}
\revised{(Day Month Year)}
\accepted{(Day Month Year)}
\published{(Day Month Year)}
\end{history}

\begin{abstract}
% TODO: 实验完成后撰写摘要
This paper presents PE-PINN, a Physics-Enhanced Physics-Informed Neural
Network framework for solving the PKN model of hydraulic fracturing.
The governing Nordgren equation is solved by embedding physical priors
directly into the network architecture, including a transient tip shape
function that analytically interpolates asymptotic regimes, a causal
arrival-time constraint, and a learnable time-power exponent. Numerical
experiments demonstrate competitive accuracy against finite difference
solutions while maintaining a mesh-free formulation.
\end{abstract}

\keywords{Physics-informed neural network; PKN model; hydraulic fracturing;
Nordgren equation; moving boundary; tip singularity.}

\section{Introduction}

Hydraulic fracturing is a crucial technology for enhancing
hydrocarbon recovery from unconventional reservoirs
\cite{Economides2000}. Accurate prediction of fracture geometry ---
the spatial and temporal evolution of fracture width, length, and net
pressure --- is essential for optimizing treatment design and evaluating
well performance. Perkins and Kern and Nordgren first proposed the PKN
model to describe the evolution of fracture length and width during
hydraulic fracturing \cite{Nordgren1972}. This classical model, together
with the KGD model \cite{Ref_KGD}, established the theoretical basis for
the simulation of hydraulic fracturing. Among these models, the PKN
model remains one of the most widely adopted in industrial practice due
to its computational tractability. The PKN model assumes a fracture of
fixed height that is much larger than its vertical extent, with
plane-strain deformation in each cross-section. Under this assumption,
the local elasticity relation yields a nonlinear governing equation that
balances the temporal evolution of fracture width, the fluid leak-off
into the surrounding formation, and the longitudinal gradient of the
volumetric flow rate governed by the fluid's viscous resistance. Despite
its apparent simplicity, this equation poses significant computational
challenges due to its strong nonlinearity, moving free boundary, and
singular behavior at the propagating tip.

The numerical solution of the PKN model has been an active research
topic since Nordgren's original finite-difference formulation
\cite{Nordgren1972}. Traditional grid-based approaches include explicit
finite-difference schemes \cite{Kemp1990}, finite element methods with
tip enrichment \cite{Garikapati2017}, and spectral discretizations
\cite{Peirce2014}. However, these methods face persistent difficulties.
The fracture tip exhibits a singular gradient that tends to infinity,
requiring special tip elements, regularization techniques
\cite{Linkov2011}, or asymptotic enrichment functions
\cite{Garikapati2017} to avoid severe accuracy degradation. The
classical PKN tip asymptote follows a one-third power-law in the
viscosity-dominated regime \cite{Nordgren1972,Kemp1990}, but a more
complex transient behavior --- smoothly transitioning from a one-third
to a three-eighths power-law depending on the local leak-off intensity
--- has been recognized in recent re-examinations
\cite{Linkov2015,Detournay2016}. Furthermore, the moving free boundary
evolves nonlinearly with time, following different power-law regimes
depending on whether the propagation is dominated by fluid storage or
leak-off. Accurately tracking this evolving crack front typically
requires adaptive remeshing or front-tracking algorithms
\cite{Peirce2014}, which introduce significant computational overhead.
Additionally, a multi-scale structure arises from the thin boundary
layer near the tip to the global fracture scale, demanding dense local
grids that can make parametric studies and real-time applications
prohibitive \cite{Detournay2016}.

Physics-Informed Neural Networks (PINNs) were introduced by Raissi
\emph{et al.} \cite{Raissi2019} as a mesh-free framework that embeds
the governing partial differential equations directly into the loss
function of a neural network. PINNs exploit automatic differentiation
for exact derivative computation and can naturally incorporate sparse
observational data. Since their introduction, PINNs and their variants
--- including domain-decomposition extensions such as XPINN
\cite{Jagtap2020} and cPINN \cite{Jagtap2020b}, gradient-enhanced
formulations \cite{Yu2022}, and multi-scale architectures
\cite{Moseley2023} --- have been successfully applied across a broad
spectrum of problems in computational mechanics. Recently, PINNs have
been applied to the field of subsurface engineering, including flow and
transport in porous media \cite{Hassan2024}, history-matching in
fractured reservoirs \cite{Ali2025}, and two-phase flow in heterogeneous
formations \cite{Wang2024}. However, the application of PINNs to
hydraulic fracture propagation --- characterized by its unique
combination of a moving free boundary, singular tip asymptotics, and
multi-regime transient behavior --- remains relatively unexplored.

A small number of recent studies have begun to bridge PINNs with
hydraulic fracture modeling. \cite{Yang2023} developed a
physics-informed surrogate model for predicting hydraulic fracture
geometry using deep learning, demonstrating the feasibility of replacing
computationally expensive simulators with trained neural networks.
\cite{Tang2024} proposed a PINN formulation with moving boundary
constraints for modeling hydraulic fracturing, making an important step
toward handling the free-boundary nature of the problem. However,
existing PINN-based approaches for hydraulic fracturing still face
critical limitations. First, most existing PINN formulations for
fracture problems do not explicitly encode the asymptotic tip behavior
into the network architecture. Without such prior knowledge, the network
must learn sharp near-tip gradients purely from PDE residuals, which is
notoriously difficult for gradient-based optimization
\cite{Krishnapriyan2021}. The transient transition between the
one-third and three-eighths power-law tip asymptotics has not been
addressed in any existing PINN work. Second, the PKN governing equation
is physically valid only in the region where the fracture has already
arrived. Existing PINN formulations do not enforce this causal
constraint, leading to non-physical residuals being minimized in regions
not yet reached by the fracture front. Third, the early-time behavior
of fracture width exhibits a power-law singularity whose exponent
depends on the dominant physical regime. Learning this temporal scaling
directly from data is inefficient and limits generalization to time
ranges not covered during training.

To address these challenges, this paper proposes PE-PINN
(Physics-Enhanced PINN), a novel framework for the PKN model that
explicitly embeds physical prior knowledge into the network architecture.
The proposed method decomposes the fracture width into three physically
meaningful components: a learnable time-power term that captures the
early-time singular scaling, an analytically derived transient tip shape
function that smoothly interpolates between the one-third and
three-eighths asymptotic regimes, and a neural network that learns the
remaining smooth component. This physics-enhanced representation
explicitly separates the singular time and space dependencies, allowing
the network to learn only a well-behaved residual function. Building on
this decomposition, the arrival time of the crack tip at each spatial
location is computed by numerically inverting the crack-length evolution
relation, and the PDE residual evaluation is restricted to the physically
meaningful region where the fracture has already arrived. This causal
masking ensures
physical consistency. Additionally, the time-power exponent is
constrained to a physically plausible range and optimized jointly with
the network parameters, automatically discovering the correct temporal
scaling. A two-stage training strategy combining Adam \cite{Kingma2015}
for global exploration with L-BFGS \cite{Liu1989} for fine local
convergence is employed. In contrast to existing PINN formulations for
hydraulic fracturing \cite{Yang2023,Tang2024}, the proposed method
encodes the asymptotic tip structure analytically rather than relying on
the optimizer to discover it, which significantly reduces the learning
burden on the network. Numerical experiments demonstrate that the
proposed method achieves competitive accuracy against conventional
finite difference solutions while maintaining a mesh-free formulation.
Systematic ablation studies and sensitivity analyses further confirm
that each component of the proposed framework contributes measurably to
overall accuracy.

The remainder of this paper is organized as follows. Section~2 formulates
the PKN governing equations. Section~3 describes the proposed PE-PINN
framework, including the network design, loss function, training
strategy, and experimental configuration. Section~4 presents the
numerical results. Section~5 concludes the paper.

% ============================================================
% Section 2: Problem Formulation
% ============================================================
\section{Problem Formulation}

The PKN model describes the propagation of a hydraulic fracture of
constant height $H$ within a homogeneous, linear elastic reservoir
\cite{Nordgren1972}. The fracture is assumed to be much longer than
its height, resulting in plane-strain deformation in each vertical
cross-section perpendicular to the propagation direction. The fluid
flow within the fracture is laminar and governed by lubrication
theory. The fracture propagates in the viscosity-dominated regime,
and fluid leak-off from the fracture faces into the surrounding
formation follows Carter's one-dimensional leak-off model
\cite{Detournay2016}.

Under the plane-strain assumption, the net pressure $p(x,t)$ is
proportional to the local fracture width $W(x,t)$ at each vertical
cross-section \cite{Nordgren1972}:
\begin{equation}
  p(x,t) = \frac{E'}{4H}\,W(x,t),
  \label{eq:elasticity}
\end{equation}
where $E' = E/(1-\nu^2)$ is the plane-strain Young's modulus. The
fluid flow within the fracture is described by the lubrication
equation, which relates the local flow rate to the pressure gradient.
Combining the elasticity and lubrication equations with the mass
conservation condition that accounts for Carter's leak-off yields the
Nordgren equation:
\begin{equation}
  \frac{\partial W}{\partial t} + \frac{C_L}{\sqrt{t - \tau(x)}}
  - \frac{E'}{256 H^2\mu}\,\frac{\partial^2(W^4)}{\partial x^2}
  = 0,
  \label{eq:nordgren}
\end{equation}
where $C_L$ is the leak-off coefficient, $\mu$ is the fluid viscosity,
and $\tau(x)$ is the arrival time of the fracture tip at the spatial
coordinate $x$. The boundary condition at the wellbore ($x = 0$)
relates the injection rate to the fracture width gradient through the
flux condition. At the propagating tip ($x = L(t)$), the fracture width
vanishes.

The fracture length $L(t)$ evolves nonlinearly with time, following
distinct power-law regimes in the storage-dominated ($L \propto
t^{4/5}$) and leak-off-dominated ($L \propto t^{1/2}$) limits
\cite{Nordgren1972}. A harmonic mean interpolation smoothly connects
these asymptotic limits.

A critical feature of the PKN model is the singular behavior of the
fracture width near the propagating tip. In the classical
viscosity-dominated regime, the width follows a one-third power-law
$W \sim (1 - x/L)^{1/3}$ as $x \to L(t)$
\cite{Nordgren1972,Kemp1990}. However, a more general analysis reveals
a transient tip shape that smoothly transitions from the one-third
power-law to a three-eighths power-law depending on the local balance
between fluid storage and leak-off \cite{Linkov2015}. This transient
behavior is characterized by a dimensionless tip shape function
$\Omega_{\rm tip}(\xi)$, where $\xi = 1 - x/L(t)$ denotes the
dimensionless distance from the tip. Accurately capturing this
transition is essential for resolving the near-tip solution without
excessive mesh refinement.

% ============================================================
% Section 3: Methodology
% ============================================================
\section{Methodology}

\subsection{PE-PINN design}

Physics-Informed Neural Networks (PINNs) were introduced by Raissi
\emph{et al.} \cite{Raissi2019} as a mesh-free framework for solving
partial differential equations. A fully connected neural network
approximates the solution $u(x, t; \theta)$, and the governing PDE
residual is embedded directly into the loss function. Automatic
differentiation is employed to compute the required derivatives exactly.

The proposed PE-PINN framework addresses three fundamental challenges
that conventional PINNs face when applied to the PKN model with a
moving boundary and singular tip behavior.

\textbf{Physics-enhanced ansatz.} Rather than directly approximating
the fracture width with a generic neural network, the proposed method
decomposes the solution into three physically meaningful components.
The time-power factor captures the early-time singular scaling, the
transient tip shape function analytically encodes the asymptotic
crack-tip behavior, and the neural network learns the remaining smooth
component. This decomposition explicitly separates the singular
dependencies, allowing the network to learn only a well-behaved
residual function.

\textbf{Causal arrival-time constraint.} The PKN governing equation is
physically valid only in the region where the fracture has already
arrived. The proposed method computes the arrival time $\tau(x)$ at
each spatial location by numerically inverting the $L(t)$ relation
using Brent's root-finding method. The PDE residual is evaluated only
where $t > \tau(x)$ and $x \leq L(t)$, ensuring that the network is not
penalized for violating physics in regions not yet reached by the
fracture front.

\textbf{Learnable time-power exponent.} The early-time temporal scaling
exponent is treated as a trainable parameter constrained to a
physically plausible range. This parameter is optimized jointly with the
network weights, automatically discovering the correct scaling without
manual tuning. Its converged value provides interpretable physical
insight into the dominant flow regime.

\textbf{Network architecture.} The neural network component consists of
a fully connected feed-forward architecture with four hidden layers of
128 neurons and Tanh activation. The output is constrained to be
positive through a Softplus activation.

\subsection{Loss function and training}

The total loss function comprises the PDE residual loss and the
boundary condition loss. The PDE residual is computed only within the
causally valid region. A smooth temporal factor is introduced at the
wellbore to reconcile the initial and boundary conditions at the
origin.

Collocation points are generated using Latin Hypercube Sampling, with
additional refinement in the near-tip region to resolve the sharp
gradients. Training proceeds in two stages: the Adam optimizer with
exponential learning rate decay for global exploration, followed by the
L-BFGS optimizer with strong Wolfe line search for fine convergence.

\subsection{Experimental configuration}

The physical parameters adopted in the numerical experiments are
summarized in Table~\ref{tab:params}.

\begin{table}[t]
\tbl{Physical parameters.\label{tab:params}}
{\begin{tabular}{@{}lcl@{}} \toprule
Parameter & Symbol & Value \\ \colrule
Plane-strain Young's modulus & $E'$ & [TBD] \\
Fracture height & $H$ & [TBD] \\
Fluid viscosity & $\mu$ & [TBD] \\
Leak-off coefficient & $C_L$ & [TBD] \\
Injection rate & $Q_0$ & [TBD] \\
Simulation time & $T_{\rm max}$ & [TBD] \\ \botrule
\end{tabular}}
\end{table}

A finite difference (FD) solver is implemented as the reference
baseline. The FD scheme employs a uniform spatial grid with explicit
time stepping and a velocity-based tip propagation condition.

To quantify the contribution of each component, four model variants
are constructed for ablation analysis: the full PE-PINN,
PINN, a variant without the time-power term, a variant without the
transient tip shape function, and a plain PINN without any of the
proposed components. In addition, a systematic sensitivity analysis
is conducted by varying the number of hidden layers (depth
$D = 2, 3, 4, 5, 6$) and the number of neurons per layer (width
$W = 32, 64, 128, 256, 512$), yielding 25 configurations.

% ============================================================
% Section 4: Results and Discussion
% ============================================================
\section{Results and Discussion}

\subsection{Model validation}

% TODO: 等实验结果完成后填充
% - 与 FDM 的裂缝宽度曲线对比 (Fig. 1, 2)
% - L(t) 演化对比
% - 定量误差对比表 (Table 2, L2 error, max error)
% - 计算效率分析 (PINN vs FDM)

The proposed method is evaluated against the finite difference
baseline. Fracture width profiles at representative times and the
crack length evolution are compared.

\subsection{Ablation study}

% TODO: 等实验结果完成后填充
% - 四个变体的裂缝宽度曲线对比 (Fig. 3)
% - 定量误差对比表 (Table 3)
% - 分析: 去掉 time-power 导致早期精度下降;
%         去掉 tip shape 导致尖端精度下降;
%         去掉两者 (plain PINN) 无法收敛

The ablation study quantifies the individual contribution of each
component of the proposed framework.

\subsection{Sensitivity analysis}

% TODO: 等实验结果完成后填充
% - 热力图: D × W → L2 error (Fig. 4)
% - 分析: 更深更宽的网络提升精度但有递减;
%         中等深度和宽度达到最优性价比
% - learnable alpha 收敛值的物理可解释性讨论

The sensitivity analysis characterizes the influence of network depth
and width on prediction accuracy.

% ============================================================
% Section 5: Conclusion
% ============================================================
\section{Conclusion}

This paper presented PE-PINN, a Physics-Enhanced PINN framework for the PKN
model of hydraulic fracture propagation. The proposed method addresses
three fundamental challenges of applying PINNs to moving boundary
problems: the transient tip shape function analytically encodes the
transition between asymptotic regimes, the causal arrival-time
constraint restricts the PDE residual to the physically valid region,
and the learnable time-power exponent adaptively captures the temporal
scaling behavior. Numerical experiments demonstrate competitive
accuracy against a conventional finite difference solver while
maintaining a mesh-free formulation. Future work will extend the
framework to more complex fracture propagation scenarios.

\section*{Acknowledgments}

This work was supported by [FUNDING INFORMATION TO BE ADDED].

\section*{References}

\begin{thebibliography}{99}
\markboth{Authors' Names}{Short Running Head}

% === [E1] 工程背景 ===
\bibitem[{Economides and Nolte(2000)}]{Economides2000}
Economides, M.~J. and Nolte, K.~G. (eds.) [2000] \emph{Reservoir Stimulation},
  3rd edn. (John Wiley \& Sons, Chichester).

% === [E2-E5] PKN模型核心文献 (你现有的4篇) ===
\bibitem[{Geertsma and de Klerk(1969)}]{Ref_KGD}
Geertsma, J. and de Klerk, F. [1969] ``A rapid method of predicting width and
  extent of hydraulically induced fractures,''
  \emph{J. Pet. Technol.} \textbf{21}(12), 1571--1581.

\bibitem[{Nordgren(1972)}]{Nordgren1972}
Nordgren, R.~P. [1972] ``Propagation of a vertical hydraulic fracture,''
  \emph{Soc. Pet. Eng. J.} \textbf{12}(4), 306--314.

\bibitem[{Kemp(1990)}]{Kemp1990}
Kemp, L.~F. [1990] ``Study of Nordgren's equation of hydraulic fracturing,''
  \emph{SPE Prod. Eng.} \textbf{5}(3), 311--314.

\bibitem[{Linkov(2015)}]{Linkov2015}
Linkov, A.~M. [2015] ``A reexamination of the classical PKN model of hydraulic
  fracture,''
  \emph{Int. J. Eng. Sci.} (to appear).

\bibitem[{Detournay(2016)}]{Detournay2016}
Detournay, E. [2016] ``A review of PKN-type modeling of hydraulic fractures,''
  \emph{J. Pet. Sci. Eng.} (to appear).

% === [E3-E4] 传统数值方法 ===
\bibitem[{Garikapati \emph{et~al.}(2017)Garikapati, Verhoosel, van Brummelen
  and Diez}]{Garikapati2017}
Garikapati, H., Verhoosel, C.~V., van Brummelen, E.~H. and Diez, P. [2017]
  ``An XFEM-based numerical strategy for the PKN model of hydraulic fracture,''
  \emph{Int. J. Numer. Methods Eng.} (to appear).

\bibitem[{Peirce(2014)}]{Peirce2014}
Peirce, A. [2014] ``Modeling multi-scale processes in hydraulic fracture
  propagation using the implicit level set algorithm,''
  \emph{Comput. Methods Appl. Mech. Eng.} \textbf{283}, 881--908.

\bibitem[{Linkov(2011)}]{Linkov2011}
Linkov, A.~M. [2011] ``On efficient simulation of hydraulic fracturing in terms
  of particle velocity,''
  \emph{Int. J. Eng. Sci.} \textbf{52}, 77--88.

% === [E5] PINN基础 ===
\bibitem[{Raissi \emph{et~al.}(2019)Raissi, Perdikaris and Karniadakis}]{Raissi2019}
Raissi, M., Perdikaris, P. and Karniadakis, G.~E. [2019] ``Physics-informed
  neural networks: A deep learning framework for solving forward and inverse
  problems involving nonlinear partial differential equations,''
  \emph{J. Comput. Phys.} \textbf{378}, 686--707.

% === [E6-E9] PINN变体与发展 ===
\bibitem[{Jagtap and Karniadakis(2020)}]{Jagtap2020}
Jagtap, A.~D. and Karniadakis, G.~E. [2020] ``Extended physics-informed neural
  networks (XPINNs): A generalized space-time domain decomposition based deep
  learning framework for nonlinear partial differential equations,''
  \emph{Commun. Comput. Phys.} \textbf{28}(5), 2002--2041.

\bibitem[{Jagtap \emph{et~al.}(2020)Jagtap, Kharazmi and Karniadakis}]{Jagtap2020b}
Jagtap, A.~D., Kharazmi, E. and Karniadakis, G.~E. [2020] ``Conservative
  physics-informed neural networks on discrete domains for conservation laws:
  Applications to forward and inverse problems,''
  \emph{Comput. Methods Appl. Mech. Eng.} \textbf{365}, 113028.

\bibitem[{Yu \emph{et~al.}(2022)Yu, Lu, Meng and Karniadakis}]{Yu2022}
Yu, J., Lu, L., Meng, X. and Karniadakis, G.~E. [2022] ``Gradient-enhanced
  physics-informed neural networks for forward and inverse PDE problems,''
  \emph{Comput. Methods Appl. Mech. Eng.} \textbf{393}, 114823.

\bibitem[{Moseley \emph{et~al.}(2023)Moseley, Markham and Nissen-Meyer}]{Moseley2023}
Moseley, B., Markham, A. and Nissen-Meyer, T. [2023] ``Finite basis
  physics-informed neural networks (FBPINNs): A scalable domain decomposition
  approach for solving differential equations,''
  \emph{Adv. Neural Inf. Process. Syst.} \textbf{36}.

% === [E10-E12] PINN在地质力学/多孔介质中的应用 ===
\bibitem[{Hassan and Labak(2024)}]{Hassan2024}
Hassan, A. and Labak, J. [2024] ``Physics-informed neural networks for
  energy-based modeling of non-Newtonian fluid flow and fracture propagation
  in subsurface formations,''
  \emph{Energy} (to appear).

\bibitem[{Ali \emph{et~al.}(2025)Ali, Yan and Gunning}]{Ali2025}
Ali, A., Yan, B. and Gunning, J. [2025] ``History-matching of imbibition flow
  in multiscale fractured porous media using physics-informed neural networks
  (PINNs),''
  \emph{Comput. Methods Appl. Mech. Eng.} \textbf{437}, 117784.

\bibitem[{Wang \emph{et~al.}(2024)Wang, Guo and Wu}]{Wang2024}
Wang, S., Guo, B. and Wu, Y.-S. [2024] ``Physics-informed neural network
  simulation of two-phase flow in heterogeneous and fractured porous media,''
  \emph{Adv. Water Resour.} \textbf{189}, 104730.

% === [E13] 移动边界/Stefan问题 ===
\bibitem[{Shkolnikov \emph{et~al.}(2024)Shkolnikov, Soner and
  Tissot-Daguette}]{Shkolnikov2024}
Shkolnikov, M., Soner, H.~M. and Tissot-Daguette, V. [2024] ``Deep level-set
  method for Stefan problems,''
  \emph{J. Comput. Phys.} \textbf{502}, 112829.

% === [E14-E15] PINN+水力压裂 (你现有文献6,7) ===
\bibitem[{Yang \emph{et~al.}(2023)}]{Yang2023}
Yang, K., et al. [2023]
  ``Physics-informed surrogate modeling for hydro-fracture geometry prediction
  based on deep learning,''
  \emph{to appear}.

\bibitem[{Tang \emph{et~al.}(2024)}]{Tang2024}
Tang, S., et al. [2024]
  ``Physics-informed neural network with moving boundary constraints for
  modeling hydraulic fracturing,''
  \emph{to appear}.

% === [E16] PINN failure modes ===
\bibitem[{Krishnapriyan \emph{et~al.}(2021)Krishnapriyan, Gholami, Zhe, Kirby
  and Mahoney}]{Krishnapriyan2021}
Krishnapriyan, A.~S., Gholami, A., Zhe, S., Kirby, R.~M. and Mahoney, M.~W.
  [2021] ``Characterizing possible failure modes in physics-informed neural
  networks,''
  \emph{Adv. Neural Inf. Process. Syst.} \textbf{34}.

% === [E17-E18] 优化器 ===
\bibitem[{Kingma and Ba(2015)}]{Kingma2015}
Kingma, D.~P. and Ba, J. [2015] ``Adam: A method for stochastic optimization,''
  \emph{Proc. 3rd Int. Conf. Learn. Represent. (ICLR)} (San Diego, CA).

\bibitem[{Liu and Nocedal(1989)}]{Liu1989}
Liu, D.~C. and Nocedal, J. [1989] ``On the limited memory BFGS method for large
  scale optimization,''
  \emph{Math. Program.} \textbf{45}, 503--528.

\end{thebibliography}

\end{document} 